\section{问题描述与控制方程}

\subsection{一维 Riemann 问题}

\subsubsection{坐标系与控制方程}
对于一维问题，坐标取为以右为正方向的一维直线坐标系。
在此坐标系下，无粘可压 Euler 方程为
\begin{equation}
\left\{
\begin{aligned}
\frac{\partial \rho}{\partial t}
&+
\frac{\partial (\rho u)}{\partial x}
=0,\\[6pt]
\frac{\partial (\rho u)}{\partial t}
&+
\frac{\partial (\rho u^2+p)}{\partial x}
=0,\\[6pt]
\frac{\partial (\rho E)}{\partial t}
&+
\frac{\partial (\rho E u+p u)}{\partial x}
=0.
\end{aligned}
\right.
\end{equation}

其中，单位质量流体的总能量为
\begin{equation}
E=e+\frac{u^2}{2}
=
\frac{p}{\rho(\gamma-1)}
+
\frac{u^2}{2}.
\end{equation}

\subsubsection{左右初始状态}

下面列出的左右状态均可理解为以间断位置为原点的局部写法。实际计算中，
程序把间断放在 \(x=x_0\) 处，因此条件 \(x<0\) 和 \(x\geq0\) 分别对应
\(x<x_0\) 和 \(x\geq x_0\)。

\noindent 1. 激波间断--接触间断--膨胀波问题之 Sod 问题
\begin{equation}
\begin{cases}
x<0: & (\rho,u,p)_L=(1,\ 0,\ 1),\\
x\geq 0: & (\rho,u,p)_R=(0.125,\ 0,\ 0.1).
\end{cases}
\end{equation}

\noindent 2. 激波间断--接触间断--膨胀波问题之 Lax 问题
\begin{equation}
\begin{cases}
x<0: & (\rho,u,p)_L=(0.445,\ 0.698,\ 3.528),\\
x\geq 0: & (\rho,u,p)_R=(0.5,\ 0,\ 0.571).
\end{cases}
\end{equation}

\noindent 3. 双膨胀波问题之亚声速后退问题
\begin{equation}
\begin{cases}
x<0: & (\rho,u,p)_L=(1,\ -2,\ 4),\\
x\geq 0: & (\rho,u,p)_R=(1,\ 2,\ 4).
\end{cases}
\end{equation}

\noindent 4. 双膨胀波问题之 Sjogreen supersonic Expansion 问题
\begin{equation}
\begin{cases}
x<0: & (\rho,u,p)_L=(1,\ -2,\ 0.4),\\
x\geq 0: & (\rho,u,p)_R=(1,\ 2,\ 0.4).
\end{cases}
\end{equation}

\noindent 5. 接触间断--双膨胀波问题
\begin{equation}
\begin{cases}
x<0: & (\rho,u,p)_L=(1,\ -0.2,\ 0.5),\\
x\geq 0: & (\rho,u,p)_R=(0.5,\ 0.5,\ 0.5).
\end{cases}
\end{equation}

\noindent 6. 接触间断--双激波间断问题
\begin{equation}
\begin{cases}
x<0: & (\rho,u,p)_L=(0.4,\ 0.5,\ 1.0),\\
x\geq 0: & (\rho,u,p)_R=(1,\ -0.5,\ 0.9).
\end{cases}
\end{equation}

\noindent 7. 接触间断问题
\begin{equation}
\begin{cases}
x<0: & (\rho,u,p)_L=(10,\ 1.0,\ 2.0),\\
x\geq 0: & (\rho,u,p)_R=(1.0,\ 1.0,\ 2.0).
\end{cases}
\end{equation}

\subsubsection{计算目标}

本项目采用 MacCormack 格式求解一维 Riemann 问题的数值解，并将数值解与
半解析解（精确解）进行对比。具体研究目标如下。

\begin{enumerate}[leftmargin=2em]
    \item 求解上面列出的几个 Riemann 初值问题，具体如下：
    \begin{itemize}
        \item 将 MacCormack 数值解与精确解进行对比；
        \item 比较加入和不加入人工粘性时的计算结果；
        \item 研究人工粘性经验参数 $\beta$ 对数值解的影响；
        \item 研究 CFL 数对数值解的影响，并讨论人工粘性系数与 CFL 数之间
        并非简单单调关系的现象；
        \item 以密度作为人工粘性开关函数，并进一步比较采用压力或速度作为
        开关函数时的差异；
        \item 在不加入人工粘性的情况下，通过加密网格考察色散效应是否得到改善，
        并分析 MacCormack 格式固有色散误差与网格分辨率之间的关系。
    \end{itemize}

    \item 对于 Sjogreen supersonic Expansion 超声速后退问题，重点考察真空区：
    \begin{itemize}
        \item MacCormack 格式能否稳定计算真空或近真空状态；
        \item 真空区内密度、压力和速度应如何处理；
        \item 数值保护、人工粘性和网格分辨率对结果的影响。
    \end{itemize}

    \item 对于接触间断--双膨胀波问题，分析接触间断和两侧膨胀波的捕捉效果。

    \item 对于接触间断--双激波问题，分析激波位置、间断厚度及数值振荡。

    \item 对于纯接触间断问题，考察 MacCormack 格式对接触间断的分辨能力，

    并在条件允许时与其他高阶或高分辨率格式进行比较。通过该算例说明不同数值格式
    各自擅长捕捉的波结构存在差异，不能仅依据格式阶数判断其在所有问题上的表现。
\end{enumerate}

\subsubsection{Project 2 与 Project 0 精确解之间的关系}
Project 0 给出同一 Riemann 问题的精确解程序。虽然程序内部在求解非线性方程时仍要用迭代方法，
但它对应的是精确 Riemann 解，可以作为 Project 2 数值解的参考解。Project 2 采用 MacCormack 格式计算近似解，
再与 Project 0 的结果比较，用来检查误差、稳定性和间断位置。

最终输出和比较时主要使用原始量 \(\bm W=(\rho,u,p)^{\mathrm T}\)。
\begin{equation}
\bm{W}(x,0)=
\begin{cases}
(\rho_L,u_L,p_L)^{\mathrm T}, & x<x_0,\\
(\rho_R,u_R,p_R)^{\mathrm T}, & x>x_0.
\end{cases}
\end{equation}

\subsection{一维 Euler 方程}

\begin{equation}
\frac{\partial \bm{Q}}{\partial t}
+
\frac{\partial \bm{F}(\bm{Q})}{\partial x}
=0,
\end{equation}

\begin{equation}
\bm{Q}=
\begin{bmatrix}
\rho\\
\rho u\\
\rho E
\end{bmatrix},
\qquad
\bm{F}=
\begin{bmatrix}
\rho u\\
\rho u^2+p\\
u(\rho E+p)
\end{bmatrix}.
\end{equation}

总能量满足

\begin{equation}
E=\frac{p}{\rho(\gamma-1)}+\frac{u^2}{2}.
\end{equation}
